% ---
% filename : oibt_autorisations_20260517_628_autorisations_limitees_oibt.tex
% id : oibt_autorisations_20260517_628
% title : "Types d'autorisations d'installer limitées selon l'OIBT"
% domain : "OIBT"
% subdomain : "Autorisations"
% tags : ["oibt", "autorisations_limitées", "inspection", "dt"]
% difficulty : 1
% time_solve : 8
% has_octave : false
% author : "om"
% ---
\documentclass[a4paper,11pt,answers,addpoints]{exam}
% ---------------------------------------------------------
% LANGUE, ENCODAGE, FONTS
% ---------------------------------------------------------
\usepackage[utf8]{inputenc}
% ---------------------------------------------------------
% GÉOMÉTRIE ET MISE EN PAGE
% ---------------------------------------------------------
\usepackage[left=1.7cm,right=1.5cm,top=2cm,bottom=1cm]{geometry}
\usepackage{microtype}
\usepackage{float}
\usepackage{needspace}

% ---------------------------------------------------------
% MATHÉMATIQUES
% ---------------------------------------------------------
\usepackage{amsmath, amssymb, bm, esvect}

\newcommand{\V}{\overrightarrow}
\def\vect#1{\overrightarrow{\strut#1}}
\providecommand{\Degrees}[0]{\ensuremath{^\circ}}

% ---------------------------------------------------------
% UNITÉS ET GRANDEURS (SIunitx)
% ---------------------------------------------------------
\usepackage{siunitx}
\sisetup{output-decimal-marker={,}}
\DeclareSIUnit \var {var}
\DeclareSIUnit \VA {VA}
\DeclareSIUnit \kVA {kVA}
\DeclareSIUnit[quantity-product={}] \degree{\text{\textdegree}}

% Permet la césure dans les nombres avec unités (évite les Underfull \hbox)
%\AllowBreakAtQQ

% ---------------------------------------------------------
% GRAPHIQUES : TikZ + CircuitikZ (sans tkz-fct qui cause des erreurs spath3)
% ---------------------------------------------------------
\usepackage{tikz}
\usetikzlibrary{
	arrows.meta,
	intersections,
	decorations.pathreplacing,
	positioning
}

\usepackage[european,straightvoltages,EFvoltages]{circuitikz}
\usepackage{pgfplots}
\pgfplotsset{compat=1.12}

% Symbole étoile-triangle personnalisé
\newcommand{\wye}{%
	\mathbin{\tikz[x=1ex,y=1ex]{%
			\draw[line width=.1ex] (0,0)--(45:1)--++(-45:1) (45:1)--++(0,1);
	}}%
}

% ---------------------------------------------------------
% TABLEAUX
% ---------------------------------------------------------
\usepackage{tabularx}
\usepackage{tabu}

% ---------------------------------------------------------
% HYPERLIENS
% ---------------------------------------------------------
\usepackage[colorlinks=true,
menucolor=red,
breaklinks=true,
urlcolor=blue,
linkcolor=blue]{hyperref}

% ---------------------------------------------------------
% COMMANDES PERSONNELLES
% ---------------------------------------------------------
\newcommand\nomauteur{bdminasmorgul@protonmail.com}
\newcommand\dateTe{18.05.2026}
\newcommand\brancheTe{DT}

% ---------------------------------------------------------
% STYLE DE L'EXERCICE
% ---------------------------------------------------------
\pagestyle{headandfoot}
\firstpageheadrule
\runningheadrule

% En-tête avec largeur fixe pour éviter les dépassements
\firstpageheader{\brancheTe\ (\nomauteur)}{Élève : }{\makebox[3.5cm][r]{\dateTe\ \thepage/\numpages}}
\runningheader{\brancheTe\ (\nomauteur)}{Élève : }{\makebox[3.5cm][r]{\dateTe\ \thepage/\numpages}}

% Pied de page
\newcommand{\totalpointtts}{%
	\begin{tikzpicture}[remember picture, overlay]
		\node[anchor=south west, inner sep=0pt, outer sep=0pt]
		at ([xshift=0.5cm,yshift=0.5cm]current page.south west)
		{\rotatebox{90}{Total des points : \numpoints}};
	\end{tikzpicture}%
}

\firstpagefooter{\hspace*{\fill}\totalpointtts}{}{}
\runningfooter{\hspace*{\fill}\totalpointtts}{}{}

% Réglages exam
\printanswers
\addpoints
\pointsinmargin
\bracketedpoints

% ---------------------------------------------------------
% LISTINGS (code)
% ---------------------------------------------------------
\usepackage{xcolor}
\usepackage{listings}

\lstdefinestyle{octavestyle}{
	language=Matlab,
	basicstyle=\ttfamily\small,
	keywordstyle=\color{blue},
	commentstyle=\color{green!50!black},
	stringstyle=\color{red},
	stepnumber=1,
	frame=single,
	breaklines=true,
	breakatwhitespace=true,
	tabsize=4
}

\usepackage{enumitem}
\setlist[enumerate,1]{label=\alph*), ref=\alph*)}
\newenvironment{explication}{%
	\par\noindent\textbf{Explication. }\itshape
}{\par\normalfont}

\begin{document}
	\unframedsolutions
	\pointsinmargin
	\bracketedpoints
	
\section*{Préparation DT PQ CFC ELMO,IELE,PELE}
\begin{questions}
\question
En tant que responsable technique au sein d'une structure industrielle, vous devez vous assurer que les interventions électriques sont réalisées par du personnel qualifié et autorisé. Selon l'Ordonnance sur les installations à basse tension (OIBT), l'Inspection fédérale (ESTI) peut délivrer des autorisations limitées pour des cadres spécifiques.

\begin{parts}
	\part[3] Quels sont les trois types d'autorisations d'installer limitées que l'Inspection peut délivrer~?
	\part[3] L'autorisation pour les travaux effectués sur des installations propres à l'entreprise permet-elle contractuellement d'effectuer la modification d'installations en aval d'un coupe-surintensité d'abonné~?
	\part[3] Quel type d'organisme doit impérativement assurer le suivi technique en cours d'emploi des électriciens d'exploitation mentionnés dans l'autorisation~?
	\part[3] Une entreprise titulaire d'une autorisation de raccordement est-elle autorisée à remplacer un moteur électrique raccordé à demeure~?
\end{parts}

\begin{solution}
\begin{itemize}
	\item \textbf{Types d'autorisations} : Les trois types sont les autorisations pour installations propres à l'entreprise (Art. 13), pour installations spéciales (Art. 14) et pour le raccordement de matériels électriques (Art. 15).
	\item \textbf{Installations propres} : Oui, l'Art. 13 al. 3 let. b autorise explicitement la modification d'installations en aval d'un coupe-surintensité d'abonné ou de dispositifs de protection des circuits terminaux.
	\item \textbf{Suivi technique} :  Selon l'OIBT Art. 13 al. 4, le suivi technique doit être assuré sans interruption par un organisme d'inspection accrédité.
	\item \textbf{Raccordement} : Oui, selon l'OIBT Art. 15, l'autorisation de raccordement donne le droit de raccorder ou de remplacer les matériels électriques raccordés à demeure, ou destinés à l'être, s'ils sont mentionnés dans l'autorisation.
\end{itemize}
\end{solution}
\end{questions}
\end{document}
